\begin{ouverture}[calcul]{8}{Fractions décimales et virgule}{Partie 2 · Fractions et nombres décimaux}
{Dixièmes · Centièmes · Écriture à virgule}
\objectif{Lire et décomposer des fractions décimales.}
\objectif{Relier unités, dixièmes et centièmes.}
\objectif{Passer d’une fraction décimale à l’écriture à virgule.}
\prerequis{Fractions sur la demi-droite graduée.}
\end{ouverture}
\pagestyle{otMain}\essentiel\fontsize{11}{15}\selectfont
\phantomsection\label{cours-8-1}
\essSection{1}{Dixièmes et centièmes}
\begin{definition}Partager l’unité en 10 donne des dixièmes : $\frac1{10}$. Partager chaque dixième en 10 donne des centièmes : $\frac1{100}$. Une unité vaut 10 dixièmes ou 100 centièmes ; un dixième vaut 10 centièmes.\end{definition}
\label{fig-cours-8-1-cours-droite-dixiemes}\figureGeo{cours-droite-dixiemes}

\begin{exemple}$\frac{23}{10}=2+\frac3{10}$ : vingt-trois dixièmes remplissent 2 unités et il reste 3 dixièmes. De même, $\frac{235}{100}=2+\frac3{10}+\frac5{100}$ : on décompose rang par rang.\end{exemple}
\begin{savoirfaire}[Lire des dixièmes]Sur une demi-droite graduée en dixièmes, compter les pas après l’entier : le point situé 4 pas après 2 repère $2+\frac4{10}$, c’est-à-dire $\frac{24}{10}$.\end{savoirfaire}
\clearpage
\phantomsection\label{cours-8-2}
\essSection{2}{Le codage de la virgule}
\begin{definition}L’écriture à virgule code une somme de fractions décimales : $35{,}78$ désigne $35+\frac7{10}+\frac8{100}$. À gauche de la virgule, la partie entière ; à droite, le chiffre des dixièmes puis celui des centièmes.\end{definition}\begin{center}\renewcommand{\arraystretch}{1.35}\begin{tabular}{c|cccc}\toprule
Nombre & dizaines & unités & dixièmes & centièmes\\\midrule
$35{,}78$ & 3 & 5 & 7 & 8\\
$6{,}09$ &  & 6 & 0 & 9\\
\bottomrule\end{tabular}\end{center}
\begin{exemple}$6{,}09=6+\frac9{100}$ : le zéro montre qu’il n’y a aucun dixième. Sans lui, on écrirait $6{,}9$, un nombre différent. On lit « 6 unités et 9 centièmes ».\end{exemple}
\begin{attention}$0{,}5$ et $0{,}05$ ne désignent pas le même nombre : 5 dixièmes pour l’un, 5 centièmes pour l’autre. En revanche $2{,}30=2{,}3$ : trente centièmes font exactement trois dixièmes.\end{attention}

\clearpage\rubrique{Découverte}
\begin{blocLivre}{ex-8-1}\begin{exercice}[1]{Des dixièmes à lire}
Écrire $\frac7{10}$ puis $\frac{23}{10}$ en toutes lettres, puis écrire $\frac{23}{10}$ avec un entier et une fraction.
\end{exercice}\end{blocLivre}
\begin{blocLivre}{ex-8-2}\begin{exercice}[1]{Sur la demi-droite}
L’unité est partagée en dixièmes. Lire les nombres repérés par A et B, en écriture fractionnaire.
\label{fig-ex-8-2}\figureGeo{dec-2}
\end{exercice}\end{blocLivre}
\begin{blocLivre}{ex-8-3}\begin{exercice}[1]{Unités, dixièmes, centièmes}
Compléter : $1=\frac{\square}{100}$ ; $\frac3{10}=\frac{\square}{100}$ ; $\frac{50}{100}=\frac{\square}{10}$.
\end{exercice}\end{blocLivre}
\begin{blocLivre}{ex-8-4}\begin{exercice}[1]{Vers la virgule}
Écrire avec une virgule : $35+\frac7{10}+\frac8{100}$, puis $\frac4{10}+\frac2{100}$.
\end{exercice}\end{blocLivre}
\begin{blocLivre}{ex-8-5}\begin{exercice}[1]{Depuis la virgule}
Décomposer $6{,}09$ puis $14{,}3$ avec des fractions décimales.
\end{exercice}\end{blocLivre}

\rubrique{Entraînement}
\begin{blocLivre}{ex-8-6}\begin{exercice}[2]{Fractions en virgule}
Écrire avec une virgule : $\frac{45}{10}$ ; $\frac{307}{100}$ ; $\frac8{100}$.
\end{exercice}\end{blocLivre}
\begin{blocLivre}{ex-8-7}\begin{exercice}[2]{Le rang de chaque chiffre}
Dans $58{,}36$, donner le chiffre des dixièmes, le chiffre des centièmes, puis la valeur du chiffre 3.
\end{exercice}\end{blocLivre}
\begin{blocLivre}{ex-8-8}\begin{exercice}[2]{Deux écritures, un nombre}
Recopier et compléter le tableau : fraction décimale et écriture à virgule de « quatre unités et sept centièmes » et de « douze dixièmes ».
\end{exercice}\end{blocLivre}
\begin{blocLivre}{ex-8-9}\begin{exercice}[2]{Avec la monnaie}
Un centime est un centième d’euro. Écrire en euros, avec une virgule : 3 euros et 45 centimes ; 275 centimes ; 8 centimes.
\end{exercice}\end{blocLivre}
\begin{blocLivre}{ex-8-10}\begin{exercice}[2]{Des zéros qui comptent}
Écrire avec une virgule : cinq centièmes, puis cinq dixièmes. Ces deux nombres sont-ils égaux ?
\end{exercice}\end{blocLivre}

\rubrique{Problèmes}
\begin{blocLivre}{ex-8-11}\begin{exercice}[3]{La grille aux cent cases}
Une grille compte 100 cases égales ; 37 sont colorées. Écrire la fraction de la grille colorée, puis ce nombre avec une virgule. Et si on colore 3 cases de plus ?
\label{fig-ex-8-11}\figureGeo{dec-13}
\end{exercice}\end{blocLivre}
\begin{blocLivre}{ex-8-12}\begin{exercice}[3]{Le compteur en dixièmes}
Un compteur affiche $45{,}8$ et augmente d’un dixième à chaque pas. Quels nombres affiche-t-il après un pas, puis après deux pas ?
\end{exercice}\end{blocLivre}
\begin{blocLivre}{ex-8-13}\begin{exercice}[3]{Les sauts mesurés}
Trois sauts mesurent $3{,}45$ m, $3{,}5$ m et $3{,}09$ m. Décomposer chaque mesure avec des fractions décimales, puis donner le chiffre des dixièmes de chacune.
\end{exercice}\end{blocLivre}
\begin{blocLivre}{ex-8-14}\begin{exercice}[3]{Fractions bien connues}
Écrire $\frac12$, $\frac14$ et $\frac34$ avec une virgule, en passant par des centièmes.
\end{exercice}\end{blocLivre}
\begin{blocLivre}{ex-8-15}\begin{exercice}[3]{Vrai ou faux ?}
Milo affirme : « $0{,}7$, c’est 7 centièmes. » Sara affirme : « $2{,}30$ et $2{,}3$ sont égaux. » Qui a raison ? Justifier avec des fractions décimales.
\end{exercice}\end{blocLivre}

\clearpage
